% Plan view of the supporting intermediate-direction rule, seen along the
% inward surface normal.
% One panel per configured offset direction. Each panel carries the surface
% tangents, the configured offset vector theta_offset, and the tangential lever
% selected by the rule in Section 2.7.3. The lever is perpendicular to
% the configured offset direction, as recorded by the right-angle marker at the origin.
%
% The configured direction is drawn as a single arrow rather than a two-headed
% axis, because the rule depends on its sense and not only on its line, and the
% arrowed arc gives the in-plane angle measured from t_1. All four offsets
% carry the same magnitude; only the direction changes between panels.
%
% Palette per FIGURE_STYLE.md: black for the frame, red for the offset, blue
% for the compliance-centre lever, matching r_c in the moment-bookkeeping
% figure. Solid lines only, no grey and no white.
\begin{tikzpicture}[
    font=\small,
    ax/.style={-{Latex[length=1.8mm]}, black},
    cmd/.style={-{Latex[length=2.2mm]}, thick, red},
    swp/.style={-{Latex[length=1.6mm]}, red},
    lev/.style={-{Latex[length=2.2mm]}, thick, blue},
    lbl/.style={font=\footnotesize, inner sep=1.8pt},
  ]

% The angle label sits at \arad. Only the -45 panel needs a larger radius: its
% arc bisector runs just under the +t_1 axis, where the t_1 label already sits.
\foreach \dx/\ua/\arc/\arad/\angtext/\name in {%
    0/0/0/0.95/{}/{about $t_1$},
    3.4/-45/1/1.25/{$-45^\circ$}/{$-45^\circ$ from $t_1$},
    6.8/45/1/0.95/{$+45^\circ$}/{$+45^\circ$ from $t_1$},
    10.2/90/1/0.95/{$+90^\circ$}/{about $t_2$}}
{
  \begin{scope}[xshift=\dx cm]
    \pgfmathsetmacro{\ra}{\ua+90}
    \pgfmathsetmacro{\amid}{\ua/2}

    % ---- configured offset, drawn in its own sense ----------------------
    % Drawn before the tangents so that in the two panels where the offset
    % lies along a tangent, the thin black axis stays visible on top of it.
    \draw[cmd] (0,0) -- (\ua:1.15);
    \node[lbl, red] at (\ua:1.42) {$\theta_{\mathrm{offset}}$};

    % ---- surface tangents ----------------------------------------------
    % The tangent labels sit below and to the left of their arrow tips, so
    % that in the panels where the offset lies along a tangent they clear the
    % theta_offset label instead of sitting on it.
    \draw[ax] (0,0) -- (0.80,0) node[lbl, below] {$t_1$};
    \draw[ax] (0,0) -- (0,0.80) node[lbl, left] {$t_2$};

    % ---- arc from t_1 to the configured direction, with its angle --------
    % The angle is set well outside the arc radius so the value reads clear of
    % the arc itself.
    \ifnum\arc=1
      \draw[swp] (0.46,0) arc (0:\ua:0.46);
      \node[lbl, red] at (\amid:\arad) {\angtext};
    \fi

    % ---- selected tangential lever --------------------------------------
    \draw[lev] (0,0) -- (\ra:1.10);
    \node[lbl, blue] at (\ra:1.34) {$r_{c,t}$};

    % ---- right angle between the two ------------------------------------
    % The marker spans the quadrant from theta_offset to r_c,t, so the scope is
    % rotated by \ua and not by \ra: rotating by the lever direction put the
    % square in the quadrant beyond r_c,t, away from the pair it marks.
    \begin{scope}[rotate=\ua]
      \draw[black, thin] (0.15,0) -- (0.15,0.15) -- (0,0.15);
    \end{scope}

    \node[lbl] at (0,-1.95) {\name};
  \end{scope}
}

\end{tikzpicture}
